\documentclass[12pt]{article}
\usepackage[margin=1in]{geometry}
\usepackage{graphicx}
\usepackage{amsmath}
\usepackage{amsfonts}
\usepackage{hyperref}
\usepackage{caption}
\usepackage{fancyhdr}
\usepackage{titlesec}
\usepackage{listings}
\usepackage{xcolor}

\title{Simulating Mechanical Energy Transfer in Rowing}
\author{Gabriel Cirker}
\date{\today}

\pagestyle{fancy}
\fancyhf{}
\rhead{Gabriel Cirker}
\lhead{Simulating Mechanical Energy Transfer in Rowing}
\rfoot{\thepage}
\setlength{\headheight}{14.5pt}

\begin{document}

\sloppy

\maketitle

\begin{abstract}
Rowing shells are propelled by intermittent strokes of multiple rowers. During each drive phase, rowers apply force through the oar blades to generate thrust on the boat, whereas during recovery phases, the boat coasts and decelerates under hydrodynamic drag. We present a simplified physics-based simulation of an eight-oared shell (8+) that models the boat and crew as a single mass subject to periodic thrust pulses and quadratic water drag. The model reproduces key features of rowing dynamics, including the rapid acceleration at stroke initiation and the oscillatory steady-state speed attained under constant power input.

Over a 20-second race-pace simulation, the shell reaches an average speed of about 5.8~m$\cdot$s$^{-1}$ with approximately 6–7\% of the rowers’ work converted into kinetic energy of the system, while the remaining 93–94\% is dissipated through drag. We explore how performance varies with drag coefficient, total mass, and stroke rate. The results align with on-water observations: reducing drag or increasing stroke rate raises the steady speed~\cite{rowinginmotion2016, hofmijster2010}, whereas adding mass mainly affects acceleration with minimal impact on cruise speed.

The model’s simplicity makes it analytically tractable and highlights the dominance of drag losses in rowing. We discuss its limitations (e.g., the lack of intra-stroke rower motion and fluid memory effects) and potential refinements, including the incorporation of rower movements and asynchronous crew coordination. Overall, this work provides a clear quantitative illustration of how rowers’ mechanical energy input partitions into sound kinetic energy versus losses, reinforcing the critical importance of drag minimization and effective technique in the sport of rowing.
\end{abstract}



\section{Introduction}

Rowing a shell presents a complex interplay between human power and boat mechanics. During a race, an eight (8+) must overcome significant resistance from the water and air. The primary opposing force is hydrodynamic drag, which increases roughly with the square of boat velocity~\cite{pulman2007}. As a result, achieving higher speeds requires more power dramatically: in fact, doubling the speed demands about an eight-fold increase in power output~\cite{dudhia2000}. Rowers intuitively sense this well-known physical reality and have been documented in rowing physics literature~\cite{dudhia2000, tucklazauskas1996}. 

\newpage

Drag in rowing has multiple components—frictional (skin drag along the hull), form drag, wave drag, and minor aerodynamic drag from the hull and rowers~\cite{pulman2007, baudouinhawkins2002}. Among these, hydrodynamic (water) drag is typically dominant at racing speeds~\cite{baudouinhawkins2002}. For instance, Sanderson and Martindale (1986) estimated that hydrodynamic resistance constitutes the bulk of energy loss in rowing, far outweighing aerodynamic effects under most conditions~\cite{sandersonmartindale1986}. This means that most of the athletes’ power is expended in overcoming water resistance, with only a small fraction translating into the boat’s kinetic energy or useful motion~\cite{rowinginmotion2016}.

Because races are often won or lost by mere seconds, understanding how rowers’ power is converted (or lost) is crucial for optimizing performance~\cite{baudouinhawkins2002, kleshnevnewsletter}. Prior studies have highlighted two broad strategies to maximize speed: (1) increasing the adequate propulsive power, and (2) minimizing losses to drag and other inefficiencies~\cite{sandersonmartindale1986, hofmijster2010}. The first involves improving the rowers’ power output and blade work, while the second requires technique adjustments to reduce speed fluctuations and drag. Notably, maintaining a more constant velocity is energetically favorable since any oscillation around a given average speed incurs extra drag penalties~\cite{rowinginmotion2016, sandersonmartindale1986}. In an ideal scenario with zero velocity variation, the energy spent against drag is minimal for a given average speed~\cite{rowinginmotion2016, sandersonmartindale1986}. 

In practice, however, the cyclic nature of the stroke—with alternating drive (acceleration) and recovery (deceleration) phases—inevitably causes the boat’s speed to oscillate each stroke. This phenomenon is known colloquially as “checking” the boat during the recovery~\cite{rowinginmotion2016, pulman2007}. Researchers and coaches have long recognized that excessive check (speed drop) between strokes wastes energy and should be minimized~\cite{smithloschner2002, hofmijster2010}. Techniques such as slower, smoother recoveries and coordinated crew motion are taught to mitigate these losses~\cite{hofmijster2010, smithloschner2002}. Sanderson and Martindale’s classic work suggested that optimizing the rowers’ motion on the recovery (e.g., controlling how quickly the crew moves back toward the catch) can reduce the velocity dip and improve efficiency~\cite{sandersonmartindale1986, smithloschner2002}. Similarly, Smith and Loschner (2002) found that lengthening the recovery (i.e., spending more time on the slide) decreases peak deceleration and reduces drag power losses~\cite{smithloschner2002}. In essence, smoother coordination that keeps the shell speed as steady as possible is beneficial for performance.

Beyond technique, there has also been interest in alternative crew coordination schemes. A conventional assumption is that all rowers should drive in synchrony (in-phase) for optimal propulsion. However, a counterintuitive hypothesis is that slight alternation (out-of-phase or “antiphase” rowing) could yield a more constant overall thrust, thereby reducing speed fluctuations~\cite{debrouwer2013, debrouwer2013}. Theoretically, if half the crew pulls while the others recover in a staggered pattern, the boat might experience more continuous propulsion with less check. De Brouwer et al. (2013) tested this idea on linked ergometers, finding that perfectly alternating strokes indeed led to more minor velocity fluctuations and less power lost to cyclic decelerations~\cite{debrouwer2013, cuijpers2015}. In their experiments, crews rowing out-of-phase transferred more of their power into useful work on the system compared to fully synchronous crews~\cite{debrouwer2013, cuijpers2015}. This suggests potential performance gains from antiphase coordination. 

However, they also noted drawbacks: antiphase rowing was less stable and harder to execute precisely, and real boats face balance and timing challenges that make asynchronous rowing impractical under race conditions~\cite{debrouwer2013, cuijpers2015}. Indeed, attempts at “syncopated” rowing in actual 8+ boats (where rowers slightly offset their drives) have historically not been widely adopted, likely due to these stability and timing issues~\cite{rowinginmotion2016}. Thus, nearly all crews continue to row in-phase, accepting some efficiency loss in exchange for maintainable rhythm and balance. Our model assumes synchronous rowing, as is standard practice, but we revisit the implications of this choice in the Discussion.

In parallel with biomechanical studies, physics-based models of rowing have been developed to analyze rowing mechanics quantitatively. Early analytical models and empirical studies established the fundamental equations of motion for a rowing shell~\cite{sandersonmartindale1986, baudouinhawkins2002}. The fundamental dynamic equation is simply Newton’s second law applied to the mass of the boat-plus-rowers system: the net force (propulsive minus drag) equals mass times acceleration~\cite{pulman2007}. Various researchers have constructed simulation models to capture aspects of the stroke. For example, Brearley and de Mestre (1996) created a semi-analytic model of the rowing stroke to investigate efficiency gains~\cite{brearleydemestre1996}. Lazauskas (1997) developed a performance prediction model for rowing races, introducing numerical simulations for shell dynamics over a race course~\cite{lazauskas1997}. His work and others’ identified different possible force-time profiles for rowers’ power application and suggested that a more symmetrically distributed effort through the stroke could be energetically optimal~\cite{lazauskas1997}. 

Van Holst (2004) developed a simulation that models the rower’s mass moving on the slide, enabling deeper analysis of how intra-stroke crew motion affects the boat’s speed oscillation~\cite{vanholst2004}. More recently, Mola (2014) presented a high-fidelity multi-physics model that coupled rowers’ movements with fluid dynamic simulations (including computational fluid dynamics for water flow)~\cite{mola2014, mola2014}. Such advanced models can capture effects like fluid memory, hull flex, and detailed blade-water interactions, at the cost of high complexity and computational expense~\cite{mola2014}. 

In contrast, the model I use here is deliberately simplified to focus on the core energy transfer mechanisms. We lump the rowers and boat into a single rigid system and represent the thrust as an idealized periodic force. Despite ignoring many fine details (e.g., the rowers’ center-of-mass motions, blade hydrodynamics, and wave generation), a simple model can still reveal the primary determinants of performance in rowing~\cite{pulman2007, rowinginmotion2016}. By isolating the fundamental physics—a cyclic input of work against a velocity-dependent drag—we can clearly illustrate how much of the rowers’ work goes into kinetic energy versus is lost to drag, and how changing physical parameters alters the outcome. The following sections describe the model setup and equations, present simulation results for an 8+ shell under various conditions, and discuss the insights and limitations of this simplified approach in the context of existing rowing research.

\newpage

\section{Methods}

\subsection{Equations of Motion}

We consider a rowing shell and crew of total mass $m$ moving in one dimension (along the boat’s direction of travel). The external forces acting on the system are:
\begin{enumerate}
    \item the thrust $F_{\text{thrust}}(t)$ produced by the rowers during the drive portion of each stroke, directed forward, and
    \item the drag force $F_d(v)$ resisting motion, directed opposite to the velocity $v(t)$.
\end{enumerate}

Applying Newton’s second law, the governing equation of motion is
\begin{equation}
    m \frac{dv}{dt} = F_{\text{thrust}}(t) - F_d(v). 
    \label{eq:eom}
\end{equation}

The resistive force is modeled using a standard quadratic drag law~\cite{pulman2007, dudhia2000}:
\begin{equation}
    F_d(v) = k_d\, v^2 .
    \label{eq:draglaw}
\end{equation}

Here, $k_d$ is an effective drag coefficient that encapsulates hydrodynamic drag (and a small contribution from aerodynamic drag) for the 8+ at race speeds. Quadratic dependence on $v$ is typical for a body moving through a fluid at moderate Reynolds numbers and has been shown to approximate rowing shell drag measurements well~\cite{buckmannharris2014, baudouinhawkins2002}. In our case, $k_d$ may be interpreted as
\[
    k_d = \frac{1}{2}\rho C_d A,
\]
where $\rho$ is the water density and $C_d A$ is an effective drag area of the hull. Rather than measure it directly, we will infer $k_d$ such that the model’s steady-state speed matches empirical expectations for an 8+ crew (as described in Section~3).

\newpage

\subsection{Thrust Model}

Rowers apply power in a cyclic fashion. During each stroke, the \emph{drive phase} is when the oar blades are in the water, and the rowers push, generating propulsion. In the subsequent \emph{recovery phase}, the blades are out of the water as rowers return to the catch, and no propulsive force acts on the boat. We model the thrust as an idealized square wave alternating with each stroke~\cite{lazauskas1997, brearleydemestre1996}. 

We define a stroke period $T$ (seconds per stroke, reciprocal of stroke rate) and a drive duration $\delta T$, where $0 < \delta < 1$ is the duty fraction of the stroke spent driving. Over each stroke cycle, the thrust is

\begin{equation}
F_{\text{thrust}}(t) =
\begin{cases}
F_{\text{drive}}, & 0 \le \phi < \delta, \\[6pt]
0, & \delta \le \phi < 1, 
\end{cases}
\label{eq:thrust}
\end{equation}

where,

\[
\phi = \frac{t \bmod T}{T}
\]
is the phase of the stroke cycle. In this formulation, $F_{\text{drive}}$ is a constant peak force applied during the drive phase.

For the simulations, we choose $T = 2.0\ \text{s}$, corresponding to a stroke rate of 30 strokes per minute (a typical racing cadence), and $\delta = 0.45$ (i.e., 45\% of the stroke spent on the drive)~\cite{kleshnevnewsletter, kleshnevnewsletter}. This gives a drive duration of $0.9$ s and a recovery of $1.1$ s per stroke. The peak force is set to $F_{\text{drive}} = 3000\,\text{N}$, roughly equivalent to 375 N per rower in an eight~\cite{baudouinhawkins2002}. A 3000 N force applied over 0.9 s at 5--6 m\,s$^{-1}$ boat speed corresponds to an instantaneous power of approximately 15--18 kW (around 2 kW per athlete), consistent with peak power outputs measured for elite rowers in short bursts~\cite{hofmijster2010, kleshnevnewsletter}.

By construction, this thrust model produces a repeated ``sawtooth'' velocity profile: the shell accelerates when $F_{\text{thrust}} = F_{\text{drive}}$ during drives, then decelerates under drag alone during recoveries. Figure~\ref{fig:speed_profile} in the Results section illustrates this pattern.

It should be noted that the chosen thrust profile is highly idealized. Real rowers do not apply a perfectly rectangular force profile—the force on the oar varies throughout the drive, often peaking mid-stroke and tapering near the finish~\cite{kleshnevnewsletter, sandersonmartindale1986}. However, the square-wave approximation captures the essential feature: a nonzero average propulsive force during drives and zero during recoveries. Prior modeling efforts have employed similar simplifications for analytical tractability~\cite{brearleydemestre1996}. The duty cycle $\delta = 0.45$ is a reasonable approximation for a racing eight; a slightly longer recovery phase allows rowers to set up for the next catch. The effect of varying $\delta$ (stroke ratio) could be explored, but here we keep it fixed and later vary only the overall stroke rate $1/T$.

\newpage

\subsection{Energy Accounting}

To quantify where the rowers’ input energy is going, we track three energy terms during the simulation~\cite{lazauskas1997, vanholst2004}:

\begin{enumerate}
    \item \textbf{Work done by rowers ($E_{\text{in}}$):}  
    This is the mechanical work input to the system via the thrust force. Over some time interval,
    \begin{equation}
        E_{\text{in}} = \int F_{\text{thrust}}(t)\, v(t)\, dt ,
        \label{eq:Ein}
    \end{equation}
    the time-integral of power (force times velocity) provided by the rowers~\cite{dudhia2000}.
    
    \item \textbf{Energy dissipated by drag ($E_{\text{drag}}$):}  
    This is the work done against water (and air) through drag:
    \begin{equation}
        E_{\text{drag}} = \int F_d(v)\, v(t)\, dt 
        = \int k_d\, v(t)^3 \, dt ,
        \label{eq:Edrag}
    \end{equation}
    representing energy lost from the system to the water as heat, vortices, and other dissipative mechanisms~\cite{baudouinhawkins2002}.
    
    \item \textbf{Kinetic energy change ($\Delta E_k$):}  
    This is the change in kinetic energy of the shell+crew system:
    \begin{equation}
        \Delta E_k = \tfrac{1}{2}m \big[v(t_f)^2 - v(t_i)^2\big],
        \label{eq:DEk}
    \end{equation}
    i.e., the increase in the system’s kinetic energy over a chosen interval~\cite{pulman2007}.  
    If the simulation begins from rest, then at any time $t$, we simply have 
    \[
    \Delta E_k(t) = \tfrac{1}{2} m v(t)^2 .
    \]
\end{enumerate}

By the work–energy principle, energy conservation for our system gives
\begin{equation}
    E_{\text{in}} = E_{\text{drag}} + \Delta E_k ,
    \label{eq:energy_balance}
\end{equation}
assuming no other losses.

We then define an efficiency metric \texorpdfstring{$\eta$}{eta} as the fraction of rower input work that is converted into kinetic energy of the boat–crew system~\cite{hofmijster2010}:
\begin{equation}
    \eta = \frac{\Delta E_k}{E_{\text{in}}}\, .
    \tag{4}
    \label{eq:efficiency}
\end{equation}

In steady state (constant average speed), $\Delta E_k$ over a full stroke cycle is approximately zero, since the kinetic energy at the start and end of the cycle match. Thus, instantaneous efficiencies are more meaningful during transient acceleration. Nonetheless, the metric \texorpdfstring{$\eta$}{eta} remains useful for comparing scenarios—in our simulations, we will observe values on the order of only a few percent, reflecting that the vast majority of energy supplied by the rowers is dissipated by drag~\cite{rowinginmotion2016}.

\newpage

\subsection{Numerical Simulation}

The model equations are integrated numerically using a simple time-stepping scheme. We employ an explicit Euler integration with a fixed small time step $\Delta t$~\cite{lazauskas1997, brearleydemestre1996, vanholst2004}. At each step, the current velocity $v$ is used to evaluate the drag force $F_d = k_d v^2$. The phase of the stroke $\phi$ (based on the current time $t$) determines whether the thrust is set to $F_{\text{drive}}$ or $0$. The acceleration
\[
a = \frac{F_{\text{thrust}} - F_d}{m}
\]
is then applied to update the velocity:
\[
v \leftarrow v + a\,\Delta t .
\]

Simultaneously, we accumulate the input and dissipated energies $E_{\text{in}}$ and $E_{\text{drag}}$ by numerically integrating $F_{\text{thrust}}\, v$ and $F_d\, v$ over time, i.e., summing $Fv\,\Delta t$ at each step.

Unless otherwise specified, we simulate a $20$-second interval, corresponding to roughly 10 stroke cycles at 30 strokes per minute. This duration is sufficient for the 8+ to reach a near steady-state speed in the model. The boat starts from rest $(v = 0)$ at $t = 0$, marking the first drive. Table~\ref{tab:baseline_parameters} summarizes the baseline parameters used in the simulation.

\begin{table}[!htbp]
\centering
\caption{Baseline simulation parameters for the 8+ model.  
Mass and drag coefficient are chosen to produce realistic race-pace behavior, while thrust and stroke timing mimic a crew rowing at approximately 30 spm.}
\label{tab:baseline_parameters}
\begin{tabular}{l c}
\hline
\textbf{Parameter} & \textbf{Value} \\
\hline
Total mass, $m$ (boat + crew) & $750\ \text{kg}$ \\
Drag coefficient, $k_d$ & $40\ \text{N}\cdot\text{s}^2/\text{m}^2$ \\
Stroke period, $T$ & $2.0\ \text{s}$ (30 spm) \\
Drive duty fraction, $\delta$ & $0.45$ \\
Drive force, $F_{\text{drive}}$ & $3000\ \text{N}$ \\
Time step, $\Delta t$ & $0.001\ \text{s}$ \\
Simulation duration & $20\ \text{s}$ \\
\hline
\end{tabular}
\end{table}

\newpage
\section{Results}

\subsection{Speed Profile and Efficiency in Baseline Scenario}

Starting from rest, the model eight accelerates rapidly during the first few strokes and approaches a repeatable oscillatory speed cycle by around $t \approx 10$~s. Figure~\ref{fig:speed_profile} shows the boat’s velocity $v(t)$ over a 20~s simulation for the baseline parameters. The characteristic sawtooth pattern is evident: during each drive phase (when the crew is applying force), velocity rises quickly and peaks at the end of the drive; during the recovery phase (no thrust), the boat slows due to drag~\cite{pulman2007, baudouinhawkins2002}. After approximately 8-10 strokes, these peaks and troughs settle into a steady oscillation, indicating that the system has reached a limit-cycle speed regime. In this state, the energy input per stroke is balanced by the energy dissipated per stroke, resulting in a constant average speed.

\begin{figure}[!htbp]
    \centering
    \includegraphics[width=0.65\linewidth]{speed_plot.png}
    \caption{Simulated boat velocity $v(t)$ over a 20~s interval for the baseline 8+ configuration, showing the characteristic oscillatory ``sawtooth'' speed profile as the shell accelerates during the drive phase and decelerates during the recovery phase.}
    \label{fig:speed_profile}
\end{figure}

In the baseline case, the average speed converges to approximately $5.8\ \text{m}\,\text{s}^{-1}$ (about 11.3 knots, or a 2000~m time of roughly 5:45)~\cite{worldrowing2012}. Instantaneously, speed fluctuates between about $5.0\ \text{m}\,\text{s}^{-1}$ (end of recovery) and $6.8\ \text{m}\,\text{s}^{-1}$ (end of drive) each stroke~\cite{worldrowing2012, rowinginmotion2016}. These values lie within expected ranges for competitive eights: race-pace velocities typically fall between $5.5$ and $6.3\ \text{m}\,\text{s}^{-1}$ depending on conditions~\cite{worldrowing2012}, and a fluctuation amplitude of roughly $\pm 0.8\ \text{m}\,\text{s}^{-1}$ is consistent with measured speed oscillations in on-water data. The oscillation magnitude of about 14--15\% of the mean matches field observations showing that boat speed varies by 10--20\% over the stroke cycle in an eight~\cite{rowinginmotion2016}.

\newpage

In addition to velocity, we track the cumulative energy budget. By 20~s, the total work done by the rowers is approximately $1.4 \times 10^{5}\ \text{J}$ (140~kJ)~\cite{rowinginmotion2016, pulman2007}. Of this, around $1.33 \times 10^{5}\ \text{J}$ (133~kJ) is dissipated through drag acting on the hull. Only about $8.8 \times 10^{3}\ \text{J}$ (8.8~kJ) contributes to increasing the kinetic energy of the system as the boat accelerates from rest. The efficiency \texorpdfstring{$\eta$}{eta}, defined in Eq.~(4) as the fraction of input work converted to kinetic energy, is approximately 6.2\% over the 20~s interval. This efficiency is largest during the initial strokes when the boat accelerates most rapidly (since $\Delta E_k$ is accumulated early), and it plateaus once steady speed is achieved, at which point nearly all additional work is dissipated by drag rather than increasing kinetic energy. The final value of approximately 6\% is consistent with prior biomechanical estimates of rowing efficiency. In effect, about 94\% of the rowers’ work is absorbed by drag, with only $\sim$6\% increasing the boat’s kinetic energy. This confirms the physical reality that most of the energy expended in rowing goes into overcoming water resistance rather than accelerating the shell~\cite{rowinginmotion2016, baudouinhawkins2002}.

To visualize the thrust input, Figure~\ref{fig:thrust_profile} plots the applied force $F_{\text{thrust}}(t)$ for the baseline simulation. As designed, it appears as a periodic square wave: a constant 3000~N during each 0.9~s drive, followed by 0~N during each 1.1~s recovery~\cite{brearleydemestre1996, brearleydemestre1996}. The idealized sharp transitions produce corresponding sharp changes in acceleration in the velocity trace (Figure~\ref{fig:speed_profile}). During each drive, the net force is initially large and positive, giving rapid acceleration. As velocity increases, drag $F_d = k_d v^2$ grows until it nearly balances the drive force at the end of the drive. When the drive ends and $F_{\text{thrust}}$ drops to zero, the net force becomes negative (equal to $-F_d$), and the boat decelerates for the entire recovery. In steady state, the peak deceleration during recovery is nearly equal in magnitude to the peak acceleration during the drive, reflecting equilibrium of energy input and dissipation over each stroke.

\begin{figure}[!htbp]
    \centering
    \includegraphics[width=0.65\linewidth]{thrust_plot.png}
    \caption{Square-wave thrust input $F_{\text{thrust}}(t)$ used in the baseline simulation, alternating between 3000~N during the 0.9~s drive phase and 0~N during the 1.1~s recovery phase.}
    \label{fig:thrust_profile}
\end{figure}

\newpage

From an energy perspective, by the end of each drive, the rowers have added an increment of energy $\Delta E_{\text{in}}$ to the system. Over the same interval, the kinetic energy rises by $\Delta E_k$, while the remainder, $\Delta E_{\text{in}} - \Delta E_k$, is dissipated by drag. During the recovery, no new energy is added ($F_{\text{thrust}} = 0$), but drag continues to remove energy, reducing $\Delta E_k$. If the boat returns to the same speed at the start and end of the stroke, then all the energy added during the drive has ultimately been dissipated. Thus, in the limit-cycle regime, essentially all of the rowers’ work in a given stroke is lost to drag by the time the stroke ends, explaining why the boat does not continue to accelerate indefinitely. Increasing the steady speed requires either increasing power input (pulling harder or rating higher) or decreasing drag (e.g., improved hull design or smoother water).

Overall, the baseline simulation captures the characteristic behavior of a rowing shell: an initial acceleration phase followed by an oscillatory steady speed maintained by the balance between propulsive work and drag losses~\cite{rowinginmotion2016}. The numerical predictions—average speed near $5.8\ \text{m}\,\text{s}^{-1}$ and efficiency near 6--7\%—are plausible for an eight. The model reinforces that only a small fraction of the rowers’ mechanical work appears as kinetic energy, with the remainder continually dissipated by drag, consistent with the well-known dominance of hydrodynamic resistance in the energetics of rowing~\cite{tucklazauskas1996}.


\newpage

\subsection{Effect of Drag Coefficient (\texorpdfstring{$k_d$)}{kd}}

Drag has an intuitively large influence on rowing performance. We varied the drag coefficient $k_d$ in the model to see its impact on steady-state speed and efficiency. Figure~\ref{fig:drag_effect} summarizes the effect of altering $k_d$ by $\pm 20\%$ from the baseline (i.e.\ testing $k_d = 32$, 40, and 48~N$\cdot$s$^2$/m$^2$). A lower drag coefficient (which could correspond to a more streamlined hull or less dense fluid) allows the shell to attain a higher steady velocity under the same power input~\cite{tucklazauskas1996, day2011}. 

\begin{figure}[!htbp]
    \centering
    \includegraphics[width=0.75\linewidth]{drag_effect.png}
    \caption{Effect of varying the drag coefficient \texorpdfstring{$k_d$}{kd} on steady-state velocity and efficiency.  
    The simulation compares $k_d = 32$, $40$, and $48$~N$\cdot$s$^2$/m$^2$, illustrating how lower drag yields higher equilibrium speed and higher energy efficiency for the same power input.}
    \label{fig:drag_effect}
\end{figure}

In our runs, reducing $k_d$ by 20\% (to 32) raised the eventual average speed to about 6.3~m$\cdot$s$^{-1}$ (an increase of roughly 0.5~m$\cdot$s$^{-1}$ over baseline). Conversely, increasing $k_d$ by 20\% (to 48) dropped the steady speed to around 5.3~m$\cdot$s$^{-1}$. This aligns with the expectation that drag forces at a given speed scale with $k_d$, so a higher drag slows the boat more strongly until equilibrium is reached at a lower speed (where input power again equals drag power). 

We also observed that higher drag not only reduces top speed but also slightly lowers the efficiency \texorpdfstring{$\eta$}{eta}, because more of the rowers’ work is being dissipated rather than stored as kinetic energy. For instance, at $k_d = 48$ the efficiency over 20~s was around 5\%, compared to 6.2\% at $k_d = 40$. At $k_d = 32$, \texorpdfstring{$\eta$}{eta} was about 7\%. Thus, improvements that lower the drag (better hull design, smoother water) can yield tangible benefits in both speed and the fraction of energy translated into motion. This is consistent with engineering efforts to design low-drag racing shells~\cite{tucklazauskas1996} and with athletes’ intuitive appreciation that fast conditions (low wind, flat water) make a boat noticeably easier to row at speed.



\subsection{Effect of Total Mass \texorpdfstring{($m$)}{(m)}}

We next examined the role of the total mass $m$ (which, in rowing, can vary due to heavier or lighter crews or differences in boat equipment weight). Interestingly, changing the mass had little effect on the ultimate steady-state speed in our model. This is because in steady state the net force averaged over a stroke is zero, so drag equals the average thrust, independent of mass. In other words, at equilibrium speed, the power in equals the power out, and mass does not appear in that balance. 

Our simulations confirmed that an 8+ weighing 10\% more (825~kg vs.\ 750~kg) reached virtually the same final average speed ($\sim 5.8\ \text{m}\cdot\text{s}^{-1}$) as a boat 10\% lighter (675~kg), assuming the same power input. However, mass did noticeably affect the acceleration and transient response. A lighter boat accelerates more quickly for a given force (since $a = F/m$). In the model runs, the 675~kg shell reached 5~m$\cdot$s$^{-1}$ in fewer strokes and exhibited larger stroke-to-stroke speed fluctuations (because there is less inertia to smooth out the cyclic acceleration). The heavier 825~kg shell accelerated more slowly off the start and had smaller fractional speed oscillations (more inertia dampening the sawtooth peaks). Figure~\ref{fig:mass_effect} illustrates that mass primarily influences how rapidly the shell responds to the thrust, but the asymptotic speed (determined by continuous power vs.\ drag) remains the same. 

\begin{figure}[!htbp]
    \centering
    \includegraphics[width=0.75\linewidth]{mass_effect.png}
    \caption{Effect of varying total system mass on transient acceleration and intra-stroke velocity fluctuations.  
    Although lighter shells accelerate more quickly and exhibit larger oscillations, all masses converge to the same steady-state speed when power input is held constant.}
    \label{fig:mass_effect}
\end{figure}

\newpage

This result aligns with basic physics: adding weight to a boat does not change the power required to sustain a given speed (aside from secondary effects such as increased hull wetting area potentially changing drag, which is not included in our simple model). It does, however, change the pace at which a crew can accelerate the shell up to speed and how much the boat slows down on recoveries. Coaches recognize this effect—for example, coxswains (whose weight is dead mass) are kept as light as possible mainly to improve acceleration and turning, not because a heavier cox would make a crew slower once at full speed (neglecting minor drag differences). 

Our model’s implication that mass has little impact on steady-state speed holds true only if the crew can maintain the same power; in reality, a heavier crew might produce more absolute power, complicating direct comparisons. But holding power constant, the finding is that mass differences (on the order of tens of kilograms) will not significantly change the top speed of an eight, a point also noted in the rowing literature~\cite{celentano1974, kerr1998}.


\subsection{Effect of Stroke Rate (Cadence)}

Stroke rate is a crucial variable in racing and rowing. Higher ratings generally increase power output, but can also introduce fatigue and slight losses in per-stroke efficiency. In the model, we can isolate the mechanical effect of stroke rate by varying $T$ (the stroke period) while keeping the drive fraction $\delta$ and the drive force $F_{\text{drive}}$ constant. Essentially, at higher stroke rates, the rowers are applying the same force pulses more frequently, which means higher overall power input. For example, going from 30~spm (2.0~s per stroke) to 36~spm ($\sim1.667$~s per stroke) increases the number of drives (and thus total work done) per unit time by 20\%.

In our simulation, increasing stroke rate indeed led to a higher average speed, as expected from the greater power. At 36~spm (with 0.9~s drives and 0.767~s recovery to maintain $\delta = 0.45$), the model 8+ reached an average speed of about 6.3--6.4~m$\cdot$s$^{-1}$ given the same 3000~N drive force—roughly a 9--10\% speed increase for a 20\% higher stroke frequency. The efficiency \texorpdfstring{$\eta$}{eta} showed a slight decline at the higher rate: more of the marginal work at 36~spm went into drag, since the drag force grows nonlinearly with speed and the shell spends less time per stroke in the relatively efficient acceleration phase and more in overcoming an already high velocity.

Conversely, rowing at a lower rate of 24~spm (2.5~s per stroke) with the same force pulses resulted in a lower sustainable speed ($\sim5.0$~m$\cdot$s$^{-1}$) but a marginally higher stroke-by-stroke efficiency, because each stroke’s work was used partly to accelerate from a slower baseline each time (where drag was slightly less). These findings align with the practical observation that higher stroke rates yield higher boat speed, albeit with diminishing returns and greater energy expenditure~\cite{sandersonmartindale1986, smithloschner2002}. 

Figure~\ref{fig:stroke_rate} shows that at 36~spm the velocity curve oscillates about a higher mean and with less relative amplitude (since there is less time for the boat to slow down on each shorter recovery)~\cite{celentano1974, dudhia2000}. The model thus captures the trade-off: a higher stroke rate with the same force per stroke increases power and speed, but also incurs more total drag work (because the boat is moving faster on average, and drag power scales with \texorpdfstring{$v^3$}{v^3}~\cite{tucklazauskas1996, day2011}). 

\newpage

\begin{figure}[!htbp]
    \centering
    \includegraphics[width=0.88\linewidth]{speed_plot.png}
    \caption{Simulated velocity profiles for different stroke rates (24, 30, and 36 spm).  
    Higher ratings produce higher mean speed and reduced relative oscillation amplitude, but at the cost of substantially increased drag power.}
    \label{fig:stroke_rate}
\end{figure}

In reality, rowers usually cannot simply maintain the same force at a much higher rate without fatigue; there is an optimum cadence for a given effort and race distance. But purely mechanically, our results illustrate why sprinting at a high rating can boost speed for short periods—more strokes mean more energy input—and also why it is inefficient to rate too high if the crew cannot sustain effective force, as wasted energy fighting drag can spike. Coaches often speak of finding the ``sweet spot’’ where the boat is moving fast, but the crew is still applying effective work per stroke; our model’s quantitative view supports this qualitative balance between rate, power, and drag losses~\cite{sandersonmartindale1986, kleshnevnewsletter}.


\newpage

\section{Discussion}

Despite its simplicity, the model reproduces several key features of rowing dynamics observed in practice. First, it generates an oscillatory speed profile characteristic of the stroke cycle, with magnitudes comparable to measured boat speed fluctuations~\cite{kleshnevnewsletter}. The initial rapid acceleration of the shell at the catch and the subsequent deceleration during recovery – often termed the “check” – are clearly present. The fact that our idealized model (with an instantaneous, constant force on drives and no force on recoveries) produces a similar oscillation suggests that the root cause of speed fluctuations is fundamentally the intermittent nature of rowing propulsion combined with drag, rather than in finer details of blade hydrodynamics. This corroborates the understanding that as long as there is a recovery phase with no propulsion, speed fluctuations are inevitable~\cite{dudhia2000, smithloschner2002}. Real crews mitigate (but cannot eliminate) this by smoothing the recovery motion to minimize how much speed is lost before the next drive~\cite{smithloschner2002, smithloschner2002}. Our model’s strict on/off thrust exaggerates the sharpness of transitions, but the overall pattern holds.

Second, the model underscores the dominance of drag in the energy balance of rowing. We found efficiency values around 6\%, meaning ~94\% of the work was dissipated by drag. This aligns with rough estimates by others that, on the order of 90–95\% of metabolic energy in rowing is expended overcoming drag, with only a small fraction contributing to the kinetic or potential energy of the system~\cite{hofmijster2010}. For example, Hofmijster et al. (2010) reported that in a single scull, only ~8\% of the power generated was transferred into kinetic energy of the boat and rower, the rest going into drag and moving water~\cite{hofmijster2010, hofmijster2010}. Our 6\% for an eight is in the same ballpark, perhaps slightly lower since larger boats typically have slightly better drag efficiency due to less speed fluctuation amplitude (as indicated by the Rowing in Motion analysis showing drag efficiency ~90–98\% depending on boat class and stroke rate~\cite{sandersonmartindale1986, rowinginmotion2016}). 


\begin{figure}[!htbp]
\centering
\includegraphics[width=0.61\linewidth]{efficiency_contour.png}
\caption{Model-predicted efficiency landscape showing how efficiency \texorpdfstring{$\eta$}{eta} varies with drag coefficient $k_d$ and stroke rate. Efficiency remains low overall due to the dominance of drag, but systematic trends align with theoretical expectations.}
\label{fig:efficiency_contour}
\end{figure}

The implication is clear: rowing fast is primarily about managing and reducing drag losses, since that’s where most energy goes. Improving technique to reduce unnecessary drag (for instance, by avoiding abrupt motions that increase water resistance, or by keeping the hull trim and avoiding extra surface area in the water) can pay dividends. Likewise, technological improvements (lighter boats, smoother hull finishes, etc.) serve to shave down $k_d$ and thereby reduce energy waste. The physics also explains why rowing with even a slight tailwind or current feels easier – effectively $k_d$ is reduced, or some drag work is offset, raising the drag efficiency.

Our parametric studies provide insight into how different factors affect performance. Reducing the drag coefficient $k_d$ was the most impactful for increasing speed per given power~\cite{lazauskas1997}. This quantitatively supports boat designers’ efforts in hull shaping for minimal drag~\cite{tucklazauskas1996}. The mass effect, by contrast, primarily affected acceleration rather than top speed. This is one reason why in rowing races, the start and sprint phases (where acceleration matters) can see differences due to crew weight, while mid-race base speed is more about power and drag.

The stroke rate analysis highlights the nonlinear nature of drag power. Since drag power scales roughly as \texorpdfstring{$v^3$}{v^3}~\cite{tucklazauskas1996}, rowing at a higher speed incurs disproportionately higher drag losses. Our model demonstrated that increasing stroke rate (hence power) yields diminishing returns in speed – a form of “diminishing marginal boat speed” that rowers know well. It also showed that higher rates reduce the relative speed fluctuation, which could slightly improve drag efficiency per stroke~\cite{sandersonmartindale1986}. 


\subsection{Limitations}

Naturally, a model of this simplicity has many limitations. It treats the entire 8+ and rowers as a point mass with a single velocity. In reality, during the recovery, the rowers’ bodies move relative to the boat, and momentum exchange causes the hull to surge forward – a phenomenon known as “boat run” where the boat actually speeds up when rowers slide toward the stern~\cite{dudhia2000, pulman2007}. Our model cannot capture this internal momentum transfer because it lacks separate masses for rowers and the boat. More advanced models include the rower mass moving on the slide and have shown that this can significantly affect the velocity profile and energy distribution~\cite{vanholst2004, mola2014}. If rowers slam up the slide too quickly, they can cause extra “stern check,” increasing energy loss~\cite{smithloschner2002}.

We also do not model the interactions between the oar blade and water. Work is done on accelerating water by the blades, creating vortices and a wake, and only a fraction of the rower's power contributes to useful propulsion. Kleshnev reports blade propulsive efficiencies of 70–80\%~\cite{kleshnevnewsletter, kleshnevnewsletter}. 

\newpage

\begin{figure}[!htbp]
\centering
\includegraphics[width=0.75\linewidth]{blade_comparison_plot.png}
\caption{Comparison of hypothetical blade efficiency curves illustrating losses due to blade slip and vortex generation. Real blades convert only a portion of the handle power into boat propulsion, motivating refined blade-hydrodynamic modeling.}
\label{fig:blade_plot}
\end{figure}

By calibrating our $k_d$ to match overall performance, blade losses are effectively folded into the drag coefficient.

Another limitation is the assumption of a constant drive force of 3000~N. Real force–time profiles are more bell-shaped~\cite{kleshnevnewsletter}. Although the area under the curve (work per stroke) is most important, instantaneous acceleration and peak speeds would differ slightly under a realistic force curve.

We also neglect added mass, wave drag, and turbulence. Wave-making resistance can matter at high speeds~\cite{tucklazauskas1996, day2011}. Aerodynamic drag can also influence performance in headwinds~\cite{pulman2007}. 

Lastly, our model assumes perfect synchronization among rowers. Studies demonstrate that slight desynchronization can reduce velocity variance and drag losses~\cite{hofmijster2010, debrouwer2013}. Alternating-stroke patterns can increase system efficiency, though they are impractical on water due to stability issues~\cite{debrouwer2013}.

Despite these limitations, the simplified model provides quantitative insight into why drag dominates rowing energetics and why smooth recoveries matter so much~\cite{rowinginmotion2016}. Coaches and athletes can use such models to reason about power, drag, and performance without requiring full CFD simulations.

\newpage

\section{Conclusion}

We presented a simple physics-based simulation of rowing, treating the rowers and boat as a single mass propelled by periodic thrust pulses against a quadratic drag. Despite its simplicity, the model produces realistic boat speed profiles and offers quantitative estimates of energy usage and efficiency. In a simulated 8+ race piece, on the order of only 5–7\% of the rowers’ mechanical work appeared as kinetic energy gain of the boat/crew system, while the remaining 93–95\% was dissipated in overcoming drag. This finding reinforces the fundamental challenge in rowing: nearly all of the athletes’ power is swallowed by water resistance. The model also illustrates how changes in key parameters affect performance: lowering drag (through hull design or conditions) significantly increases speed~\cite{tucklazauskas1996}, extra mass mainly impacts acceleration but not steady speed, and higher stroke rates yield faster speeds at the cost of proportionally greater energy expenditure against drag~\cite{hofmijster2010}. These outcomes align with both empirical data and the prevailing understanding among rowers and coaches.

While the simplified model cannot capture all nuances (such as the intricacies of oar-water interaction, crew synchronization, or hull flex), it serves as a clear demonstration of the physics governing rowing speed. The fact that such a minimal model aligns well with observed behavior suggests that our fundamental representation – a cyclic forcing of a drag-limited system – is sound. Future work could refine the model by incorporating a time-varying thrust profile that more closely mimics real-world force application, adding separate masses for rowers to capture the effects of their movements on shell velocity, or coupling the simulation with fluid dynamics to improve drag characterization~\cite{mola2014, day2011}. Such extensions would enable exploration of questions such as how much “checking” the boat truly costs in terms of time over a race, or how an optimal coxswain steering strategy (minimizing yaw drag) might save energy.

In summary, this study highlights that rowing fast is largely about managing losses. By quantifying how rower power is partitioned, we gain a deeper appreciation for technical improvements that can make small but crucial differences (for example, improvements that raise efficiency from 6\% to 7\% could translate to meaningful seconds saved over 2000 m). The simplified model serves as an educational tool for rowers and coaches to understand why aspects such as streamlining, consistent rhythm, and recovery technique matter so much. It also provides a baseline for more sophisticated simulations of rowing, which can build on this framework to include additional real-world factors. Ultimately, combining physics-based models with athlete training can help push the envelope of rowing performance in a sport where every drop of energy counts.

\newpage

\appendix
\appendix
\section*{Appendix A: Baseline Simulation Code \newline(rowing\_model.py)}

To reproduce the baseline results and the plots in Figures~\ref{fig:speed_profile} and~\ref{fig:thrust_profile}, we use the following Python script. It integrates the equations of motion with an explicit Euler scheme and saves the velocity time series and thrust profile as PNG figures.

\lstset{basicstyle=\ttfamily\small, breaklines=true, frame=single}

\begin{lstlisting}[language=Python,caption={Baseline 8+ simulation used for the velocity and thrust plots.},label={lst:rowing_model}]
"""
Simulation of Mechanical Energy Transfer in Rowing
--------------------------------------------------

This script implements a simple dynamic model of an eight-person racing shell
propelled by rowers. The model balances thrust generated during the drive
phase of each stroke against hydrodynamic drag on the hull.

The governing equation for the boat's forward motion is

    m dv/dt = F_thrust(t) - k_d v^2,

with thrust modeled as a square wave in time: constant during the drive
fraction of each stroke and zero during recovery. Mechanical energy input
from the rowers and energy dissipated by drag are accumulated over time so
that an overall transfer efficiency can be estimated.

Running this script produces:
  * speed_plot.png   : boat velocity v(t) over 20 s for the baseline case
  * thrust_plot.png  : square-wave thrust profile F_thrust(t)
  * rowing_results_summary.txt : text summary of key quantities
"""

import numpy as np
import matplotlib.pyplot as plt


def simulate_rowing(
    m=750.0,
    k_d=40.0,
    T=2.0,
    delta=0.45,
    F_drive=3000.0,
    total_time=20.0,
    dt=0.001,
):
    """
    Simulate the forward motion of an 8+ shell over a given time horizon.

    Parameters
    ----------
    m : float
        Total mass of system (boat + crew) in kg.
    k_d : float
        Quadratic drag coefficient (N s^2 / m^2).
    T : float
        Stroke period in seconds (e.g., 2.0 s for 30 spm).
    delta : float
        Drive fraction of the stroke (0 < delta < 1).
    F_drive : float
        Constant thrust during the drive phase (N).
    total_time : float
        Total simulation time in seconds.
    dt : float
        Time step for explicit Euler integration (s).

    Returns
    -------
    time : ndarray
        Time array of length n_steps.
    velocity : ndarray
        Boat velocity v(t) in m/s.
    E_in : float
        Total mechanical work done by rowers (J).
    E_drag : float
        Total energy dissipated by drag (J).
    E_kin : float
        Change in kinetic energy of the system (J).
    efficiency : float
        Mechanical efficiency, E_kin / E_in (dimensionless).
    """

    n_steps = int(total_time / dt)
    time = np.zeros(n_steps)
    velocity = np.zeros(n_steps)

    # Energies
    E_in = 0.0
    E_drag = 0.0

    v = 0.0
    t = 0.0

    for i in range(n_steps):
        time[i] = t
        velocity[i] = v

        # Determine thrust from stroke phase
        phase = (t % T) / T
        if phase < delta:
            F_thrust = F_drive
        else:
            F_thrust = 0.0

        # Quadratic drag
        F_d = k_d * v**2

        # Net acceleration
        a = (F_thrust - F_d) / m

        # Power contributions for this time step
        E_in += F_thrust * v * dt
        E_drag += F_d * v * dt

        # Update state (explicit Euler)
        v = v + a * dt
        t = t + dt

    # Change in kinetic energy from initial rest
    E_kin = 0.5 * m * (velocity[-1]**2 - velocity[0]**2)
    efficiency = E_kin / E_in if E_in > 0.0 else 0.0

    return time, velocity, E_in, E_drag, E_kin, efficiency


def run_baseline_demo():
    """
    Run the baseline simulation and produce the main figures for the paper.
    """

    # Baseline parameters (match Table 1)
    m = 750.0
    k_d = 40.0
    T = 2.0
    delta = 0.45
    F_drive = 3000.0
    total_time = 20.0
    dt = 0.001

    time, velocity, E_in, E_drag, E_kin, efficiency = simulate_rowing(
        m=m,
        k_d=k_d,
        T=T,
        delta=delta,
        F_drive=F_drive,
        total_time=total_time,
        dt=dt,
    )

    # Basic summary
    avg_vel = np.mean(velocity[int(len(velocity) / 2) :])
    max_vel = np.max(velocity)

    print("Baseline rowing simulation (8+)")
    print("--------------------------------")
    print(f"Average velocity (last half): {avg_vel:.3f} m/s")
    print(f"Maximum velocity:            {max_vel:.3f} m/s")
    print(f"Total energy input E_in:     {E_in:.1f} J")
    print(f"Energy lost to drag E_drag:  {E_drag:.1f} J")
    print(f"Change in KE E_kin:          {E_kin:.1f} J")
    print(f"Mechanical efficiency:       {efficiency * 100:.2f} %")

    # Plot 1: velocity time series (speed_plot.png)
    plt.figure(figsize=(7.0, 4.0))
    plt.plot(time, velocity, linewidth=1.0)
    plt.xlabel("Time (s)")
    plt.ylabel("Boat speed v (m/s)")
    plt.title("Velocity time series for baseline 8+ simulation")
    plt.tight_layout()
    plt.savefig("speed_plot.png", dpi=300)
    plt.close()

    # Plot 2: thrust time series (thrust_plot.png)
    thrust = np.zeros_like(time)
    for i, t in enumerate(time):
        phase = (t % T) / T
        thrust[i] = F_drive if phase < delta else 0.0

    plt.figure(figsize=(7.0, 4.0))
    plt.step(time, thrust, where="post", linewidth=1.0)
    plt.xlabel("Time (s)")
    plt.ylabel("Thrust F_thrust (N)")
    plt.title("Square-wave thrust profile for baseline simulation")
    plt.tight_layout()
    plt.savefig("thrust_plot.png", dpi=300)
    plt.close()

    # Also write a small text summary file
    with open("rowing_results_summary.txt", "w") as f:
        f.write("Rowing simulation results\n")
        f.write("==========================\n")
        f.write(f"Total simulation time: {time[-1]:.2f} s\n")
        f.write(f"Average velocity (last half): {avg_vel:.3f} m/s\n")
        f.write(f"Maximum velocity: {max_vel:.3f} m/s\n")
        f.write(f"Total energy input E_in: {E_in:.1f} J\n")
        f.write(f"Energy lost to drag E_drag: {E_drag:.1f} J\n")
        f.write(f"Change in kinetic energy E_kin: {E_kin:.1f} J\n")
        f.write(f"Mechanical efficiency: {efficiency * 100:.2f} %\n")


def main():
    run_baseline_demo()


if __name__ == "__main__":
    main()
\end{lstlisting}

\newpage

\section*{Appendix B: Advanced Analysis Scripts \newline(rowing\_advanced\_models.py)}

To generate the additional figures showing the effect of drag coefficient, total mass, and stroke rate, along with the efficiency contour and illustrative blade-efficiency curves, we use a separate script:

\begin{lstlisting}[language=Python,caption={Advanced parameter sweeps and visualization used for Figures related to drag, mass, efficiency, and blade losses.},label={lst:rowing_advanced_models}]
"""
Advanced analysis for the rowing energy-transfer model
------------------------------------------------------

This script builds on the baseline simulator defined in rowing_model.py.
It provides helper routines to explore how the model behaves when key
parameters (drag coefficient, total mass, stroke rate) are varied, and
to visualize trends in efficiency.

Running this script generates the following figures:

  * drag_effect.png            : effect of drag coefficient k_d on speed
                                and efficiency
  * mass_effect.png            : effect of total mass on transient speed
  * efficiency_contour.png     : efficiency as a function of k_d and
                                stroke rate
  * blade_comparison_plot.png  : illustrative blade efficiency curves
"""

import numpy as np
import matplotlib.pyplot as plt

from rowing_model import simulate_rowing


def explore_drag_variation():
    """
    Vary the drag coefficient k_d and examine the change in
    steady-state speed and efficiency.
    """

    k_values = [32.0, 40.0, 48.0]
    avg_speeds = []
    efficiencies = []

    for k_d in k_values:
        time, velocity, E_in, E_drag, E_kin, eta = simulate_rowing(k_d=k_d)
        avg_speed = np.mean(velocity[int(len(velocity) / 2) :])
        avg_speeds.append(avg_speed)
        efficiencies.append(eta * 100.0)

    fig, ax1 = plt.subplots(figsize=(7.0, 4.0))

    ax1.plot(k_values, avg_speeds, "o-", linewidth=1.2)
    ax1.set_xlabel("Drag coefficient k_d (N s^2 / m^2)")
    ax1.set_ylabel("Average speed (m/s)")
    ax1.set_title("Effect of drag coefficient on speed and efficiency")

    ax2 = ax1.twinx()
    ax2.plot(k_values, efficiencies, "s--", linewidth=1.2)
    ax2.set_ylabel("Mechanical efficiency (%)")

    fig.tight_layout()
    fig.savefig("drag_effect.png", dpi=300)
    plt.close(fig)


def explore_mass_variation():
    """
    Vary the total mass m and examine how quickly the shell
    accelerates and how large the stroke-to-stroke fluctuations are.
    """

    masses = [675.0, 750.0, 825.0]
    total_time = 20.0
    dt = 0.001

    fig, ax = plt.subplots(figsize=(7.0, 4.0))

    for m in masses:
        time, velocity, _, _, _, _ = simulate_rowing(
            m=m, total_time=total_time, dt=dt
        )
        label = f"m = {m:.0f} kg"
        ax.plot(time, velocity, linewidth=1.0, label=label)

    ax.set_xlabel("Time (s)")
    ax.set_ylabel("Boat speed v (m/s)")
    ax.set_title("Effect of total mass on acceleration and speed profile")
    ax.legend()
    fig.tight_layout()
    fig.savefig("mass_effect.png", dpi=300)
    plt.close(fig)


def sweep_drag_and_rate_for_efficiency():
    """
    Compute the model's efficiency as a function of drag coefficient
    and stroke rate, and display the result as a contour plot.
    """

    # Drag coefficients and stroke rates to explore
    k_values = np.linspace(30.0, 50.0, 9)  # N s^2 / m^2
    rates = np.array([24.0, 28.0, 32.0, 36.0, 40.0])  # strokes per minute

    efficiency_grid = np.zeros((len(rates), len(k_values)))

    for i, rate in enumerate(rates):
        T = 60.0 / rate  # convert spm to stroke period
        for j, k_d in enumerate(k_values):
            _, _, E_in, E_drag, E_kin, eta = simulate_rowing(
                k_d=k_d, T=T, total_time=20.0
            )
            efficiency_grid[i, j] = eta * 100.0  # percent

    fig, ax = plt.subplots(figsize=(7.0, 4.0))
    X, Y = np.meshgrid(k_values, rates)
    contour = ax.contourf(
        X, Y, efficiency_grid, levels=15, cmap="viridis"
    )
    cbar = fig.colorbar(contour, ax=ax)
    cbar.set_label("Mechanical efficiency (%)")

    ax.set_xlabel("Drag coefficient k_d (N s^2 / m^2)")
    ax.set_ylabel("Stroke rate (spm)")
    ax.set_title("Efficiency landscape as a function of drag and rate")

    fig.tight_layout()
    fig.savefig("efficiency_contour.png", dpi=300)
    plt.close(fig)


def make_blade_efficiency_plot():
    """
    Construct simple hypothetical blade efficiency curves to illustrate
    how losses due to blade slip and vortex generation might vary with
    stroke phase or loading. This is an illustrative plot only and is
    not derived from a detailed CFD model.
    """

    phi = np.linspace(0.0, 1.0, 200)  # normalized stroke phase

    # Three hypothetical "blade types" with slightly different behavior
    eta_A = 0.78 - 0.10 * (phi - 0.5) ** 2
    eta_B = 0.72 - 0.18 * (phi - 0.4) ** 2
    eta_C = 0.75 - 0.14 * (phi - 0.6) ** 2

    fig, ax = plt.subplots(figsize=(7.0, 4.0))
    ax.plot(phi, eta_A * 100.0, label="Blade A")
    ax.plot(phi, eta_B * 100.0, label="Blade B")
    ax.plot(phi, eta_C * 100.0, label="Blade C")

    ax.set_xlabel("Normalized stroke phase")
    ax.set_ylabel("Blade propulsive efficiency (%)")
    ax.set_title("Illustrative blade efficiency curves")
    ax.legend(loc="best")

    fig.tight_layout()
    fig.savefig("blade_comparison_plot.png", dpi=300)
    plt.close(fig)


def main():
    explore_drag_variation()
    explore_mass_variation()
    sweep_drag_and_rate_for_efficiency()
    make_blade_efficiency_plot()


if __name__ == "__main__":
    main()
\end{lstlisting}



\begin{thebibliography}{99}

\bibitem{pulman2007}
Pulman, C. \textit{The Physics of Rowing}. University of Cambridge Technical Report, 2007.
Discussion of forces on an eight, types of drag, and basic rowing dynamics.

\bibitem{lazauskas1997}
Lazauskas, L. \textit{A Performance Prediction Model for Rowing Races}. University of Adelaide Report, 1997.
Early simulation of rowing shell speed and power requirements.

\bibitem{vanholst2004}
Van Holst, M. \textit{Simulation of Rowing}. Online report, 2004.
Introduces rower mass motion in a rowing model and its effects.

\bibitem{tucklazauskas1996}
Tuck, E.O. \& Lazauskas, L. \textit{Low Drag Rowing Shells}. Technical Report, 1996.
Analysis of hull shapes aimed at minimizing drag for racing shells.

\bibitem{sandersonmartindale1986}
Sanderson, B. \& Martindale, W. Towards optimizing rowing technique.
\textit{Med. Sci. Sports Exerc.} 18(4):454–468, 1986.
Classic study quantifying power losses in rowing and emphasizing reduced boat velocity fluctuations.

\bibitem{kleshnevnewsletter}
Kleshnev, V. \textit{Rowing Biomechanics Newsletter}, various issues (2001–2010).
Reports on rowing force curves, efficiency, and technique (e.g., boat acceleration and effectiveness).

\bibitem{baudouinhawkins2002}
Baudouin, A. \& Hawkins, D.
Investigation of biomechanical factors affecting rowing performance.
\textit{J. Sports Sci.}, 2002.
Reviews rowing mechanics; emphasizes that system velocity is dictated by the balance of propulsive and drag forces.

\bibitem{day2011}
Day, A. H., Campbell, I., Clelland, D., Doctors, L.J. \& Cichowicz, J.
Realistic evaluation of hull performance for rowing shells in unsteady flow.
\textit{J. Sports Sci.} 29(10):1059–1069, 2011.
Studies unsteady hydrodynamics of rowing hulls, including added resistance due to speed fluctuations.

\bibitem{brearleydemestre1996}
Brearley, M. \& de Mestre, N.
Modelling the rowing stroke and increasing its efficiency.
In: \textit{Proc. 3rd Int. Conf. on Mathematics and Computers in Sport}, 1996.
Early mathematical model of the stroke; suggests a symmetric force profile as energetically optimal.

\bibitem{celentano1974}
Celentano, F. et al.
\textit{Dynamics of Rowing: Mechanics and Motion of Rowers in Boat}, 1974.
Identified momentum transfer during recovery causing boat surge (classical mechanics analysis in rowing).

\bibitem{dudhia2000}
Dudhia, A. \textit{Basic Physics of Rowing}. Oxford University, 2000.
Online resource explaining rowing physics; source of the “eight times the power to double speed” example.

\bibitem{hofmijster2010}
Hofmijster, M. et al.
Effect of stroke rate on the distribution of net mechanical power in rowing.
\textit{Scand. J. Med. Sci. Sports} 20(4):e88–e95, 2010.
Shows that a notable fraction of power is lost to within-cycle velocity fluctuations in synchronized rowing.

\bibitem{debrouwer2013}
de Brouwer, A.J., de Poel, H.J. \& Hofmijster, M.
Don’t Rock the Boat: How Antiphase Crew Coordination Affects Rowing.
\textit{PLoS ONE} 8(1):e54996, 2013.
Experiment showing that antiphase (alternating) rowing can reduce power loss to velocity fluctuations and improve useful power transfer.

\bibitem{rowinginmotion2016}
Volker, F. (Rowing in Motion).
\textit{Drag Efficiency in Rowing} (blog article), June 2016.
Explains the concept of drag efficiency; notes that crews typically achieve 89–98\% drag efficiency and that a 5\% velocity variation can waste about 5\% of energy.

\bibitem{buckmannharris2014}
Buckmann, J. \& Harris, S.
An experimental determination of the drag coefficient of a men’s 8+ racing shell.
\textit{SpringerPlus} 3:512, 2014.
Field measurements of drag on an 8+; provides an experimental drag coefficient and discusses performance implications.

\bibitem{kerr1998}
Kerr, S. \textit{Balance of Racing Rowing Boats}. Hydrodynamics report, 1998.
Investigates stability and balance; shows that lower center of gravity improves stability, while higher moment of inertia (e.g., wider beam) increases drag.

\bibitem{worldrowing2012}
World Rowing.
\textit{Men’s Eight – Speed and Power}, 2012.
Notes that the men’s eight is the fastest boat class with a world-best time of 5:19.35 over 2000 m ($\approx$ 6.27 m/s)
); provides historical context for 8+ performance.

\bibitem{smithloschner2002}
Smith, R. \& Loschner, C.
Biomechanics of rowing: coordination and efficiency.
\textit{Sports Biomech.}, 2002.
Studied rower segment coordination; emphasized slow slide recovery to minimize velocity loss and drag.

\bibitem{fisa2005}
International Rowing Federation (FISA).
\textit{Rowing Technique Guidelines}, 2005.
Coaching guidelines stressing minimization of “check” and maintenance of hull speed during recovery.

\bibitem{mola2014}
Mola, A.
\textit{Multi-physics and Multilevel Fidelity Modeling of Olympic Rowing Boats}.
Ph.D. Thesis, Virginia Tech, 2014.
High-fidelity CFD-coupled model of rowing; demonstrates influence of rower motions on boat oscillations and energy ratios.

\bibitem{cuijpers2015}
Cuijpers, L.S., den Hartigh, R.J.R., Noordzij, M.L. \& van der Kamp, J.
Coordination dynamics in crew rowing: effects of antiphase vs. in-phase patterns.
(University of Groningen report / dissertation, c. 2015).
Analyses coordination dynamics in crew rowing; provides additional evidence on efficiency changes with alternating stroke patterns.

\bibitem{vanholst_oar2007}
Van Holst, M.
Energy efficiency of the rowing oar from catch to square-off.
(Technical report, c. 2007; University of Queensland repository).
Discusses energy efficiency of the oar–water interaction and estimates blade propulsive efficiency (on the order of 70–80\%).

\end{thebibliography}

\noindent(Additional background references are cited in-line; the items above correspond to the key external sources explicitly used in this paper.)

\end{document}
